\section{Rotación con estrategia de dos usuarios (versión~4.0.0)}

Esta sección aplica únicamente a la versión~\textbf{4.0.0} de Event Hub, desplegada sobre \textbf{Amazon Aurora PostgreSQL} con la contraseña del usuario master gestionada por Aurora en AWS Secrets Manager (\texttt{--manage-master-user-password}).

\subsection{Por qué hace falta distinguir master y aplicación}

Aurora, con rotación gestionada, mantiene un secreto del estilo \texttt{rds!cluster-\ldots} asociado al usuario master (en el laboratorio, \texttt{master}). Ese secreto se rota de forma monousuario: al completarse el cambio, la password master antigua deja de ser válida en el clúster.

La estrategia de \textbf{dos usuarios} no sustituye esa rotación del master. Sirve para las credenciales con las que \textbf{conecta Spring Boot} a Aurora, de modo que pueda existir una ventana de solape: mientras se actualiza un usuario, el otro sigue aceptando conexiones.

En resumen:

\begin{itemize}
    \item \textbf{Secreto master} (\texttt{rds!cluster-\ldots}): creado y rotado por Aurora.
    \item \textbf{Secreto de aplicación} (p.~ej.\ \texttt{/prod/aurora/eventhub-app}): contiene el usuario/password que importa Event Hub; aquí aplica la estrategia de dos usuarios.
\end{itemize}

\subsection{Idea de la estrategia}

Se crean dos usuarios de aplicación con los mismos privilegios sobre la base \texttt{eventhub}, por ejemplo \texttt{eventhub\_a} y \texttt{eventhub\_b}.

\begin{enumerate}
    \item El secreto de aplicación en \texttt{AWSCURRENT} apunta a uno de ellos (p.~ej.\ \texttt{eventhub\_a}).
    \item En la siguiente rotación se actualiza la password del \emph{otro} usuario (\texttt{eventhub\_b}) y se deja esa versión como \texttt{AWSPENDING} hasta validarla.
    \item Si la prueba es correcta, \texttt{AWSCURRENT} pasa a \texttt{eventhub\_b} y la versión anterior queda como \texttt{AWSPREVIOUS}.
    \item \texttt{eventhub\_a} \textbf{sigue siendo válido} en Aurora hasta la rotación siguiente: las instancias de Event Hub que aún no han recargado el secreto pueden seguir conectando.
\end{enumerate}

Las etiquetas de \textit{staging} del secreto de aplicación quedan así:

\begin{itemize}
    \item \texttt{AWSCURRENT}: credencial que deben usar las aplicaciones ahora.
    \item \texttt{AWSPENDING}: credencial nueva mientras se aplica y se prueba la rotación.
    \item \texttt{AWSPREVIOUS}: credencial anterior; con dos usuarios, suele seguir siendo aceptada por la base de datos durante el solape.
\end{itemize}

El usuario master (secreto gestionado por Aurora) se usa solo como cuenta privilegiada para crear o alterar los usuarios de aplicación; \textbf{no} es el usuario del \texttt{DataSource} de Spring Boot en este diseño.

\subsection{Comandos CLI: secreto master gestionado por Aurora}

Tras crear el clúster con \texttt{aurora.sh} (o el equivalente con \texttt{--manage-master-user-password}), se puede consultar el secreto master asociado:

\begin{lstlisting}[language=bash,basicstyle=\ttfamily\footnotesize]
DB_CLUSTER_ID=eventhub-cluster-aurora

SECRET_ARN=$(aws rds describe-db-clusters \
  --db-cluster-identifier "$DB_CLUSTER_ID" \
  --query 'DBClusters[0].MasterUserSecret.SecretArn' \
  --output text)

aws secretsmanager describe-secret --secret-id "$SECRET_ARN"
\end{lstlisting}

Para forzar una rotación inmediata de la password master (rotación gestionada por Aurora; requiere aplicación inmediata):

\begin{lstlisting}[language=bash,basicstyle=\ttfamily\footnotesize]
aws rds modify-db-cluster \
  --db-cluster-identifier "$DB_CLUSTER_ID" \
  --rotate-master-user-password \
  --apply-immediately
\end{lstlisting}

Lectura de la versión actual del secreto master (útil para administrar usuarios de aplicación desde un entorno con acceso de red a Aurora):

\begin{lstlisting}[language=bash,basicstyle=\ttfamily\footnotesize]
aws secretsmanager get-secret-value \
  --secret-id "$SECRET_ARN" \
  --version-stage AWSCURRENT \
  --query SecretString \
  --output text
\end{lstlisting}

El JSON incluye, entre otros, \texttt{username}, \texttt{password}, \texttt{host} y \texttt{port}.

\subsection{Comandos CLI: usuarios de aplicación y secreto propio}

Con una sesión SQL autenticada como master (por ejemplo con \texttt{psql} desde la EC2 del laboratorio), se crean los dos usuarios y se les conceden privilegios sobre \texttt{eventhub}:

\begin{lstlisting}[language=bash,basicstyle=\ttfamily\footnotesize]
# Ejemplo: variables tomadas del secreto master (AWSCURRENT)
export PGHOST=...   # host del clúster
export PGPORT=5432
export PGUSER=master
export PGPASSWORD=...  # password master actual
export PGDATABASE=eventhub

psql <<'SQL'
CREATE USER eventhub_a WITH PASSWORD 'cambiar-en-el-primer-alta';
CREATE USER eventhub_b WITH PASSWORD 'cambiar-en-el-primer-alta';
GRANT CONNECT ON DATABASE eventhub TO eventhub_a, eventhub_b;
GRANT USAGE ON SCHEMA public TO eventhub_a, eventhub_b;
GRANT SELECT, INSERT, UPDATE, DELETE ON ALL TABLES IN SCHEMA public
  TO eventhub_a, eventhub_b;
ALTER DEFAULT PRIVILEGES IN SCHEMA public
  GRANT SELECT, INSERT, UPDATE, DELETE ON TABLES TO eventhub_a, eventhub_b;
SQL
\end{lstlisting}

A continuación se crea el secreto de aplicación que importará Spring Boot (versión inicial con \texttt{eventhub\_a}):

\begin{lstlisting}[language=bash,basicstyle=\ttfamily\footnotesize]
APP_SECRET_NAME=/prod/aurora/eventhub-app
CLUSTER_ENDPOINT=$(aws rds describe-db-clusters \
  --db-cluster-identifier "$DB_CLUSTER_ID" \
  --query 'DBClusters[0].Endpoint' \
  --output text)

aws secretsmanager create-secret \
  --name "$APP_SECRET_NAME" \
  --description "Credenciales de aplicacion Event Hub 4.0.0 (estrategia dos usuarios)" \
  --secret-string "{
    \"username\": \"eventhub_a\",
    \"password\": \"cambiar-en-el-primer-alta\",
    \"host\": \"${CLUSTER_ENDPOINT}\",
    \"port\": 5432,
    \"dbname\": \"eventhub\"
  }"
\end{lstlisting}

En \texttt{aurora.properties} de la versión~4.0.0, el import apuntaría a ese secreto de aplicación (no al \texttt{rds!cluster-\ldots} del master):

\begin{lstlisting}[language=properties,basicstyle=\ttfamily\footnotesize]
# Ejemplo para estrategia de dos usuarios (aplicacion).
spring.config.import=aws-secretsmanager:/prod/aurora/eventhub-app
\end{lstlisting}

\texttt{application.properties} sigue resolviendo \texttt{\$\{username\}}, \texttt{\$\{password\}}, \texttt{\$\{host\}}, \texttt{\$\{port\}} y \texttt{\$\{dbname\}} como en el laboratorio.

\subsection{Rotación manual del secreto de aplicación (solape)}

Una rotación de dos usuarios, expresada con la CLI, tiene este esquema:

\begin{enumerate}
    \item Elegir el usuario \emph{inactivo} (si \texttt{AWSCURRENT} es \texttt{eventhub\_a}, se rota \texttt{eventhub\_b}).
    \item Como master, cambiar la password de ese usuario en Aurora (\texttt{ALTER USER ... PASSWORD ...}).
    \item Publicar una versión nueva del secreto de aplicación con ese usuario y password, marcada como pendiente.
    \item Tras validar la conexión, promover esa versión a \texttt{AWSCURRENT}.
\end{enumerate}

Publicar la versión nueva (\texttt{AWSPENDING} se gestiona según el flujo de rotación; con publicación directa suele quedar como nueva \texttt{AWSCURRENT} si no se usan etapas intermedias). Ejemplo de nueva versión del secreto:

\begin{lstlisting}[language=bash,basicstyle=\ttfamily\footnotesize]
# Password nueva solo para eventhub_b; eventhub_a sigue valido en Aurora.
NEW_PASSWORD='...'  # generar de forma segura

aws secretsmanager put-secret-value \
  --secret-id "$APP_SECRET_NAME" \
  --secret-string "{
    \"username\": \"eventhub_b\",
    \"password\": \"${NEW_PASSWORD}\",
    \"host\": \"${CLUSTER_ENDPOINT}\",
    \"port\": 5432,
    \"dbname\": \"eventhub\"
  }"
\end{lstlisting}

Tras \texttt{put-secret-value}, Secrets Manager deja la versión recién escrita como \texttt{AWSCURRENT} y la anterior como \texttt{AWSPREVIOUS}. En ese momento:

\begin{itemize}
    \item Las instancias que recarguen o reinicien usarán \texttt{eventhub\_b}.
    \item Las que aún tengan en memoria la password de \texttt{eventhub\_a} pueden seguir conectando mientras ese usuario no se invalide en la siguiente rotación.
\end{itemize}

Para inspeccionar versiones y etiquetas:

\begin{lstlisting}[language=bash,basicstyle=\ttfamily\footnotesize]
aws secretsmanager list-secret-version-ids \
  --secret-id "$APP_SECRET_NAME"

aws secretsmanager get-secret-value \
  --secret-id "$APP_SECRET_NAME" \
  --version-stage AWSCURRENT

aws secretsmanager get-secret-value \
  --secret-id "$APP_SECRET_NAME" \
  --version-stage AWSPREVIOUS
\end{lstlisting}

Si hace falta mover etiquetas de forma explícita entre versiones concretas:

\begin{lstlisting}[language=bash,basicstyle=\ttfamily\footnotesize]
# VERSION_CURRENT / VERSION_PREVIOUS: VersionId obtenidos con list-secret-version-ids
aws secretsmanager update-secret-version-stage \
  --secret-id "$APP_SECRET_NAME" \
  --version-stage AWSCURRENT \
  --move-to-version-id "$VERSION_NUEVA" \
  --remove-from-version-id "$VERSION_ANTERIOR"
\end{lstlisting}

\subsection{Relación con Event Hub 4.0.0}

\begin{itemize}
    \item La rotación \textbf{gestionada por Aurora} (\texttt{modify-db-cluster --rotate-master-user-password}) sigue siendo la forma correcta de renovar el master; no debe usarse como password del \texttt{DataSource} si se adopta la estrategia de dos usuarios.
    \item El solape lo aportan \texttt{eventhub\_a} / \texttt{eventhub\_b} y el secreto \texttt{/prod/aurora/eventhub-app}.
    \item Tras promover \texttt{AWSCURRENT}, las instancias de \texttt{eventhub-eventos-springboot} deben reiniciarse o refrescar el \texttt{DataSource}; hasta entonces siguen con la credencial cacheada al arranque, que con dos usuarios puede seguir siendo válida gracias al solape.
\end{itemize}
